\subsection{EXISTENCE OF UNIVERSAL MAPPINGS}

A universal mapping problem does not necessarily have a solution (Exercise 1). However, we shall show that the following conditions imply the existence of a solution:

\((\mathbf{CU}_{\mathbf{I}})\) On every product of a family of \(\Sigma\) -sets there exists a product structure of species \(\Sigma\) (§ 2, no. 4).

\((\mathrm{CU}_{\mathrm{II}})\) Let \((\mathbf{F}_i)_{i\in \mathbf{I}}\) be a family of \(\Sigma\) -sets, and for each \(i\in I\) let \(\varphi_i\) be an \(\alpha\) -mapping of E into \(\mathbf{F}_i\) . Then the mapping \((\varphi_i)_{i\in I}\) of E into \(\prod_{i\in I}F_i\) (endowed with the product structure) is an \(\alpha\) -mapping.

A subset G of a \(\Sigma\) -set F will be said to be \(\Sigma\) -admissible if the structure on F induces a structure of species \(\Sigma\) on G ( \(\S\) 2, no. 4).

(CU \(_{III}\) ) There exists a cardinal \(\mathfrak{a}\) with the following properties: for every \(\Sigma\) -set \(\mathbf{F}\) and every \(\alpha\) -mapping \(\varphi\) of \(\mathbf{E}\) into \(\mathbf{F}\) there exists a \(\Sigma\) -admissible subset \(\mathbf{G}\) of \(\mathbf{F}\) which contains \(\varphi(\mathbf{E})\) , has cardinal \(\leqslant \mathfrak{a}\) , is such that the mapping of \(\mathbf{E}\) into \(\mathbf{G}\) having the same graph as \(\varphi\) is an \(\alpha\) -mapping, and such that any two morphisms of \(\mathbf{G}\) into a \(\Sigma\) -set, which agree on \(\varphi(\mathbf{E})\) , are equal.

CST22. If the conditions \((\mathrm{CU}_{\mathrm{I}})\) to \((\mathrm{CU}_{\mathrm{III}})\) are satisfied, then the universal mapping problem for E has a solution.

We shall first show that if there exists a pair \((\mathbf{F}_{\mathbf{E}}, \varphi_{\mathbf{E}})\) which satisfies \((\mathbf{AU}_{\mathbf{I}}^{\prime})\) , then there also exists a solution of the universal mapping problem for E. For by \((\mathrm{CU}_{\mathrm{III}})\) there exists a \(\Sigma\) -admissible subset \(F_{E}^{\prime}\) of \(F_{E}\) containing \(\varphi_{E}(E)\) , such that the mapping \(\varphi_{E}^{\prime}\) of E into \(F_{E}^{\prime}\) which has the same graph as \(\varphi_{E}\) is an \(\alpha\) -mapping, and such that any two morphisms of \(F_{E}^{\prime}\) into a \(\Sigma\) -set which agree on \(\varphi_{E}(E)\) are equal. Let j be the canonical injection of \(F_{E}^{\prime}\) into \(F_{E}\) , so that \(\varphi_{E} = j \circ \varphi_{E}^{\prime}\) . For every morphism f of \(F_{E}\) into a \(\Sigma\) -set F, \(f \circ j\) is a morphism of \(F_{E}^{\prime}\) into F, and we have \(f \circ \varphi_{E} = (f \circ j) \circ \varphi_{E}^{\prime}\) . It is therefore clear that \((\mathbf{F}_{\mathbf{E}}^{\prime}, \varphi_{\mathbf{E}}^{\prime})\) satisfies \((\mathbf{AU}_{\mathbf{I}}^{\prime})\) and \((\mathbf{AU}_{\mathbf{II}}^{\prime})\) .

It remains for us to establish the existence of a pair \((\mathrm{F}_{\mathbf{E}}, \varphi_{\mathbf{E}})\) satisfying \((\mathrm{AU}_{\mathbf{I}}^{\prime})\) . Let \(s \in \mathrm{S}(x)\) be the typical characterization of the species of structures \(\Sigma\) , and consider the subset L of \(\mathfrak{P}(\mathfrak{a}) \times \mathrm{S}(\mathfrak{a}) \times \mathfrak{P}(\mathbf{E} \times \mathfrak{a})\) consisting of all triples \(\lambda = (\mathbf{X}, \mathbf{V}, \mathbf{P})\) having the following property: ``V is a structure of species \(\Sigma\) on \(X \subset a\) , and P is the graph of an \(\alpha\) -mapping of E into X (with respect to the structure V)''

(observe that we have \(S(X) \subset S(a)\) , as is easily seen by arguing step by step on the length of the echelon construction scheme \(S\) ). For every \(\lambda = (X, V, P) \in L\) , we denote by \(X_{\lambda}\) the set \(X\) endowed with the structure \(V\) , and by \(\varphi_{\lambda}\) the mapping of \(E\) into \(X_{\lambda}\) whose graph is \(P\) .

Let \(F_{E}\) be the \(\Sigma\) -set which is the product of the \(X_{\lambda}\) (it exists by \((CU_{I})\) ), and let \(\varphi_{E}\) be the mapping \(x \to (\varphi_{\lambda}(x))\) of E into \(F_{E}\) , which is an \(\alpha\) -mapping by virtue of \((CU_{II})\) . Let us show that the pair \((F_{E}, \varphi_{E})\) satisfies \((AU_{I}')\) . Given an \(\alpha\) -mapping \(\varphi\) of E into a \(\Sigma\) -set F, let G be a subset of F which satisfies the conditions stated in \((CU_{III})\) . Let j be the canonical injection of G into F, and let \(\psi\) be the mapping of E into G which has the same graph as \(\varphi\) , so that \(\varphi = j \circ \psi\) . It follows from \((CU_{III})\) that \(\psi\) is an \(\alpha\) -mapping of E into G. Since Card \((G) \leqslant a\) , there is a subset \(G'\) of a equipotent to G. Let g be a bijection of G onto \(G'\) . If we transport by means of g the structure of species \(\Sigma\) on G, there exists by definition an element \(\lambda\) of L such that \(G'\) (endowed with the transported structure) is equal to \(X_{\lambda}\) and such that \(g \circ \psi = \varphi_{\lambda}\) . Then \(f = j \circ \overline{g}^1 \circ \mathrm{pr}_{\lambda}\) is a morphism of \(\mathbf{F}_{\mathbf{E}}\) into \(\mathbf{F}\) such that \(\varphi = f \circ \varphi_{\mathbf{E}}\) , and the proof is complete.

CST23. Let \((\mathbf{F}_{\mathbf{E}}, \varphi_{\mathbf{E}})\) be a solution of the universal mapping problem for E. Then \(\varphi_{E}\) is an injection of E into \(F_{E}\) if and only if, for each pair of distinct elements x, y of E, there exists an \(\alpha\) -mapping \(\varphi\) of E into a \(\Sigma\) -set F such that \(\varphi(x) \neq \varphi(y)\) .

Since \(\varphi_{E}\) is an \(\alpha\) -mapping, the criterion is an immediate consequence of the definitions.

In this case the \(\alpha\) -mappings are said to separate the elements of E, and we do not ordinarily make any distinction, in the terminology, between the elements of E and their images under \(\varphi_{E}\) . With this convention, if \((F_{E}, \varphi_{E})\) is a solution of a universal mapping problem, and if condition \((CU_{III})\) is satisfied, then every \(\alpha\) -mapping of E into a \(\Sigma\) -set F extends uniquely to a morphism of \(F_{E}\) into F.

